\section{Social Facts and Group Control}

\textbf{Emile Durkheim}, following the positivist school of sociology, defines the \textbf{Social Fact} as social phenomena having an existence independent of the individuals comprising the social group. As such, it has a coercive effect. A corollary is that the individuals comprising the group can change, while leaving the group behavior intact.

Positivism admits only the existence of facts, not consciousness, the material world, not the spiritual. Hence Durkheim cannot provide a theory of how the social fact came to be independent of any human beings; nor can he explain how the social fact controls the group, nor even how the group came to be constituted in the first place. How are we then to understand this from the metaphysical point of view?

First of all, to be an effective force, the social fact must be experienced as a thought in consciousness, though to be efficacious, it must not be reflective self-consciousness. That is, it is experienced in its immediacy and as unquestionable as any observation of the external world; in other words, no epistemological distinction is made between the social fact and an external fact. Julius Evola describes this type of consciousness:

\begin{quotex}
the ``I" still lives only as if in a dream: it is not yet a self-consciousness, nor an autonomous principle of action: immersed in an immediate, indistinct coalescence with nature and the world, we can say that it is not so much he who thinks, speaks, and asserts himself, as much as that various forces and impulses think, speak and assert themselves in him. He therefore is only a type of medium, a passive instrument that has his very life outside himself, and he experiences everything as grace, as spontaneity, as the immediate self-manifestation of something that transcends him.\footnote{From \emph{The Individual and the Becoming of the World.}}

\end{quotex}
In the first edition of \emph{Pagan Imperialism} intended for an Italian audience, Evola makes explicit reference to Durkheim and other similar sociologists by name. Although the names are removed in the subsequent German edition, the argument is the same. Durkheim used field studies of primitive societies to reach his conclusions, on the mistaken assumption that they represent earlier stages of evolution of more developed societies. Evola comments on the stage of consciousness of such societies, although he uses the word ``totem" in much the same sense as social fact:

\begin{quotex}
the totem is the mystical soul of the group, the clan, or the race: the individual members do not feel themselves, in their blood and in their life, as anything other than incarnations of this collective spiritual force, and possessing in themselves almost no trace of personality. 

\end{quotex}
Given Evola's thesis, then, such societies do not represent human societies at all, but something else entirely. In them, individuals are fungible, essentially indistinguishable from each other. Just as in a bale of hay, one piece of straw can be replaced with another without altering the characteristics of the bale, so in such societies the individual members may change without altering the group characteristics. Thus, they persist unchanged for centuries, unless disturbed by outside forces. Evola writes:

\begin{quotex}
If the totemic force remains at this diffuse and faceless level, so to speak, and if, consequently, there are neither leaders nor subjects, and the individual members of the group are nothing more than ``placed together" (com-posti) — then we find ourselves at the lowest level of human society, at the level which borders on the subhuman, that is, the animal kingdom: this is confirmed by the fact that the totems—the mystical souls of the clan — are often regarded at the same time as the ``spirits" of particular animal species. In addition, it is interesting that, even though the totems represent masculine figures, the composition of these societies reflects above all the telluric-matriarchal type. 

\end{quotex}
So, while Durkheim's ideas may pertain to such societies, they are not valid for a Traditional civilisation, properly so called. Thus the claim of sociology to represent a universal knowledge of social groups is falsified. A true civilisation is compose of persons, not individuals. A person, as opposed to the individual, is a conscious agent. As such, he is not replaceable by any arbitrary individual, as he has a specific place and role. The aim of such a civilisation is not to keep its members in bondage to a lower animal-like form of life, but rather to further develop the personality. Evola explains:

\begin{quotex}
A civilisation, in the true and higher sense — with reference both to individuals and peoples — arises only where the totemic level is surpassed, and where the race element, also understood mystically, is not the last resort; where, beyond blood, a force of higher, meta-biological, spiritual, and ``solar" type manifests itself, which does not lead outside of life, but determines life, transforming it, refining it, giving it a form which it initially did not have, freeing it entirely from every mixture with animal life, and opening the various paths for the realisation of the various personality types. 

\end{quotex}
This is not to say that Durkheim's ideas fail to describe man as he is in a disordered untraditional ``civilisation". Stuck in a worldview of naturalism, materialism, and rationalism, modern man is no longer able to transcend the social fact. Seeing no sense in the idea of a ``virile overcoming", he is mired in sensuality. Unable to discern any connection with a higher world, he is satisfied with an animal existence. Intellectually incapable of understanding the sacred traditional sciences, he subjects himself to the limiting ideas of a scientific positivism as a self-fulfilling prophecy. The modern techniques of mass manipulation — advertising, propaganda, media — have introduced new totems to occlude the personality of modern man. Evola offers an alternative:

\begin{quotex}
Far from exhausting itself in a naturalism — as today only the ignorance or the tendentious falsification of some people are able to present it — beyond knowing the ideals of virile overcoming and of absolute liberation, in the pagan conception, the world was a living body, suffused with secret, divine and demonic forces, with meanings and with symbols, as illustrated by that saying of Olympiodorus: ``the sensible expression of the invisible". Man lived in an organic and essential connection with the forces of the world and of the supraworld, so that it could be said, with the hermetic expression, to be ``a whole within the whole, composed of all the powers": the sense which is revealed by the aristocratic doctrine of the Atman is no different. And that conception was the basis on which, as a whole in its perfect way, the corpus of the sacred traditional sciences developed. 

\end{quotex}


\flrightit{Posted on 2010-09-06 by Cologero }

\begin{center}* * *\end{center}

\begin{footnotesize}\begin{sffamily}



\texttt{VisionsOfGlory14 on 2010-09-06 at 22:41 said: }

Durkheim's theories, then, represent a will to turn man into the Totemic, mass type described above. Belief in sociological fixes for problematic Social Facts like suicide presuppose a will to make oneself compliant with the proposed fixes. For instance, Prozac as a tool to ``cure" the growing dissatisfaction with modern society presupposes an agent's compliance with the idea that they are simply a ``depressed" person in need of the medication.

Durkheim reduces religion and matters of the spirit to the same sociological level, valuing only their communal aspects. His system is an attack on the individual's ability to rise above mere ``individuality" and become a person that acts as opposed to being acted upon.


\hfill

\texttt{Matt on 2010-09-07 at 00:45 said: }

VisionsofGlory14

I very much agree with your sentiments about what appears to be the growing desire to immediately medicate people who have a deeper angst about the existential condition than what is seen on average in a given population. However, and correct me if I'm wrong, I don't agree with (if thats what you meant) that anti-depressants are solely prescribed to make someone not have to deal with the existential and metaphysical concerns/questions of the current human condition. A lot of the time that is the case, but though science is obviously not omnipotent and omniscient, one still cannot dismiss its findings (just like how someone cannot dismiss non-scientific avenues, which is unfortunately is happening more and more). And scientific findings have shown that symptoms like depression and anxiety are sometimes caused by chemical imbalances, and the drugs are then prescribed to try and correct those imbalances. I don't really have a problem with that, its when they are prescribed when a given individual does not have chemical imbalances. Its the attitude of ``this can just be solved by medication and thats that" that I am concerned about.


\hfill

\texttt{Matt on 2010-09-07 at 00:59 said: }

* The second last sentence in my previous post should have had ``its when they are prescribed when a given individual does not clearly have a chemical imbalance that I begin to have a problem".

As for Durkhiem, his system has obvious flaws, which rightly pointed out. But with that said, I very much prefer the sociological school of thought that he is associated with (functionalism and socio-biological theory) over all the post modern, feminist, marxist-frankfurt critical theory that dominated sociology for almost all of the second half of the 20th century. Durkhiem, Comte and others at least saw society in very much the same way as Tradition does as an organic composition whose inhabitants' nature/identity are not based solely on social constructs and that much of their identity is not determined by mere society. The only difference is that their worldview does not allow for possibilities that extend beyond the physical/natural.


\hfill

\texttt{Cologero on 2010-09-07 at 08:28 said: }

It's true that Comte was a counter-revolutionary; in other words, his view was that positivism would provide an intellectually acceptable ground for traditional European society, which had been previously grounded on throne and altar. In the latter part of his life, he took a mystical turn, discerning that sociology was insufficient to explain and motivate European civilisation.

The real issue is that sociology claims to be a science; how then did critical theory and its offshoots gain such a hold in the academy? Imagine geocentrists on the astronomy faculty — not even thinkable. Yet the equivalent is acceptable in the ``soft sciences".


\hfill

\texttt{Matt on 2010-09-07 at 15:31 said: }

Cologero

That is quite interesting about Comte. I didn't know he took that turn in his worldview later on in life.

As for your comment on sociology….yeah, it does seem quite hard to argue that its somehow a science when you take that point you made into account.


\hfill

\texttt{Cologero on 2010-09-08 at 21:52 said: }

He created a religion of humanity that mimics the trappings of Catholicism without the supernatural. He posits the ``Great Being", which, as far as I can tell, collects the great thoughts of humanity. As a positivist, Comte accepts only what can be experienced, and the Great Being is experiential. Maurras has a nice introduction to Comte on this web site:

\url{http://www.gornahoor.net/library/AugusteComte.pdf}

A side note involves Solovyov. He doctoral thesis was a polemic against positivism. Yet, later on in life he came to see some positive (pun intended) value in Comte's conception of the Great Being, which Solovyov related to the Diving Sophia. I have a piece on that topic in the pipeline.


\end{sffamily}\end{footnotesize}
